\documentclass[10pt]{article}
\providecommand{\FIGDIR}{fig}
\newcommand{\DIGESTDATE}{5 September 2026}
\newcommand{\PERIODSTART}{30 August 2026}
\newcommand{\PERIODEND}{5 September 2026}
% ===================== PM Research Watch - shared preamble =====================
\usepackage[T1]{fontenc}
\usepackage[utf8]{inputenc}
\usepackage{amsmath,amssymb}
\usepackage{tgpagella}
\usepackage{tgheros}
\renewcommand{\sfdefault}{qhv}
\renewcommand{\ttdefault}{lmtt}
\usepackage{microtype}
\usepackage[english]{babel}

\usepackage[a4paper,top=17mm,bottom=17mm,left=17mm,right=17mm,
            headsep=6mm,footskip=10mm]{geometry}
\usepackage{graphicx}
\usepackage{xcolor}
\usepackage{colortbl}
\usepackage{tikz}
\usetikzlibrary{calc}
\usepackage[most]{tcolorbox}
\usepackage{titlesec}
\usepackage{fancyhdr}
\usepackage{booktabs}
\usepackage{array}
\usepackage{longtable}
\usepackage{enumitem}
\usepackage{multicol}
\usepackage{ragged2e}
\usepackage[hidelinks]{hyperref}

% ------------------------------------------------------------------ palette
\definecolor{Ink}    {HTML}{14262E}
\definecolor{Deep}   {HTML}{0F4C5C}
\definecolor{Teal}   {HTML}{1F7A8C}
\definecolor{Sky}    {HTML}{4FA3C4}
\definecolor{Sage}   {HTML}{6FA287}
\definecolor{Amber}  {HTML}{E3A008}
\definecolor{Clay}   {HTML}{B5643C}
\definecolor{Coral}  {HTML}{D1495B}
\definecolor{Violet} {HTML}{6A5B9E}
\definecolor{Slate}  {HTML}{4A6572}
\definecolor{Paper}  {HTML}{FBF9F5}
\definecolor{Mist}   {HTML}{EEEAE2}
\definecolor{Rule}   {HTML}{D9D3C9}

\pagecolor{Paper}
\color{Ink}

\hypersetup{colorlinks=true, linkcolor=Deep, urlcolor=Teal, citecolor=Teal,
            pdftitle={PM Research Watch}, pdfauthor={Automated daily digest}}

% ------------------------------------------------------------------ headings
\titleformat{\section}
  {\sffamily\bfseries\large\color{Deep}}{}{0pt}
  {}[{\vspace{1.2mm}\color{Amber}\titlerule[1.1pt]}]
\titlespacing*{\section}{0pt}{5.5mm}{2.8mm}

\titleformat{\subsection}
  {\sffamily\bfseries\normalsize\color{Teal}}{}{0pt}{}
\titlespacing*{\subsection}{0pt}{3.6mm}{1.4mm}

\setlist[itemize]{leftmargin=4.2mm, itemsep=1.1mm, topsep=1.3mm, parsep=0pt,
                  label=\textcolor{Teal}{\footnotesize$\blacktriangleright$}}

% ------------------------------------------------------------------ header/footer
\pagestyle{fancy}
\fancyhf{}
\renewcommand{\headrulewidth}{0pt}
\renewcommand{\footrulewidth}{0pt}
\fancyhead[L]{\sffamily\scriptsize\color{Slate}PM RESEARCH WATCH\;\textcolor{Rule}{\textbar}\;\DIGESTDATE}
\fancyhead[R]{\sffamily\scriptsize\color{Slate}Particulate matter monitoring \& health effects}
\fancyfoot[C]{\sffamily\scriptsize\color{Slate}\thepage}

% ------------------------------------------------------------------ components
% metric tile
\newcommand{\metric}[3]{%
  \begin{tcolorbox}[enhanced, colback=white, colframe=Rule, boxrule=0.5pt,
    arc=1.6mm, left=2.2mm, right=2.2mm, top=1.8mm, bottom=1.8mm,
    borderline west={1.7mm}{0pt}{#3}, nobeforeafter, valign=center,
    height=17mm, width=\linewidth]
    \RaggedRight
    {\sffamily\bfseries\fontsize{17}{17}\selectfont\color{#3}#1}\\[0.5mm]
    {\sffamily\fontsize{6.5}{7.4}\selectfont\color{Slate}#2}
  \end{tcolorbox}}

% pull-quote / key insight box
\newtcolorbox{keybox}[1][]{enhanced, colback=Mist, colframe=Mist,
  arc=1.6mm, left=3.2mm, right=3.2mm, top=2.4mm, bottom=2.4mm,
  borderline west={1.7mm}{0pt}{Amber}, #1}

% subtopic band
\newcommand{\band}[2]{%
  \vspace{3.2mm}\noindent
  \begin{tcolorbox}[enhanced, colback=#2, colframe=#2, arc=1.2mm,
    left=2.6mm, right=2.6mm, top=1.1mm, bottom=1.1mm, boxrule=0pt,
    width=\linewidth, nobeforeafter]
    {\sffamily\bfseries\small\color{white}#1}
  \end{tcolorbox}\vspace{1.6mm}}

% one paper entry: {short cite}{title}{journal}{doi}{body}
\newcommand{\paperentry}[5]{%
  \par\vspace{1.9mm}\noindent
  {\sffamily\bfseries\footnotesize\color{Ink}#1}\;%
  {\sffamily\fontsize{7}{8}\selectfont\color{Slate}#3}\\[0.4mm]
  {\footnotesize\itshape\color{Deep}#2}\\[0.6mm]
  {\footnotesize #4}\\[0.3mm]
  {\fontsize{6.6}{7.6}\selectfont\color{Slate}\texttt{doi:}\,\href{https://doi.org/#5}{\texttt{\detokenize{#5}}}}
  \par}

% caption style
\newcommand{\figcap}[1]{\par\vspace{1.0mm}
  {\fontsize{7.2}{8.6}\selectfont\color{Slate}\textsf{\textbf{Fig.}}\; #1\par}\vspace{1.5mm}}

\providecommand{\permil}{\textperthousand}

\setlength{\parindent}{0pt}
\setlength{\parskip}{0pt}
\renewcommand{\arraystretch}{1.18}


\begin{document}

% ==================================================================== MASTHEAD
\thispagestyle{empty}
\begin{tikzpicture}[remember picture, overlay]
  \fill[Deep] (current page.north west) rectangle
      ($(current page.north east)+(0,-67mm)$);
  \fill[Teal] ($(current page.north west)+(0,-67mm)$) rectangle
      ($(current page.north east)+(0,-68.4mm)$);
  \foreach \x/\y/\r in {22/12/0.9, 41/26/1.5, 63/9/0.7, 88/31/1.1, 112/17/1.9,
                        138/29/0.8, 159/13/1.3, 178/54/1.0, 196/21/1.6,
                        30/44/0.6, 74/57/1.2, 125/48/0.9, 168/6/0.8, 191/59/0.7,
                        52/18/1.0, 100/60/0.6, 148/20/1.1, 66/49/1.4, 205/45/0.8}
    \fill[Sky, opacity=0.30] ($(current page.north west)+(\x mm,-\y mm)$) circle (\r mm);
\end{tikzpicture}

\vspace*{4mm}
{\sffamily\fontsize{7.6}{9}\selectfont\color{Sky}%
 WEEKLY ROLLUP\;\textbullet\;PhD RESEARCH WATCH\par}
\vspace{2.2mm}
{\sffamily\bfseries\fontsize{27}{29}\selectfont\color{white}%
 Particulate Matter\par}
{\sffamily\fontsize{27}{29}\selectfont\color{Sky}%
 Week in Review\par}
\vspace{5.0mm}
{\sffamily\fontsize{8.4}{10}\selectfont\color{white}%
 \mbox{Week ending \DIGESTDATE}\;\;\textcolor{Teal}{\textbar}\;\;\mbox{\PERIODSTART{} to \PERIODEND}\;\;
 \textcolor{Teal}{\textbar}\;\;\mbox{150 records pooled from 7 daily issues}\par}
\vspace{18mm}

% -------------------------------------------------------------- metric strip
\noindent
\begin{minipage}[t]{0.190\linewidth}\metric{150}{RECORDS POOLED\\(150 UNIQUE DOIs)}{Deep}\end{minipage}\hfill
\begin{minipage}[t]{0.190\linewidth}\metric{54\%}{MEASUREMENT SIDE\\(81 RECORDS)}{Teal}\end{minipage}\hfill
\begin{minipage}[t]{0.190\linewidth}\metric{10}{SUBTOPICS\\ACTIVE}{Sky}\end{minipage}\hfill
\begin{minipage}[t]{0.190\linewidth}\metric{34}{POOLED EFFECT\\ESTIMATES}{Amber}\end{minipage}\hfill
\begin{minipage}[t]{0.190\linewidth}\metric{16}{ABSTRACT-FREE\\(11\% OF WEEK)}{Clay}\end{minipage}

\vspace{6mm}

% -------------------------------------------------------------- week in one box
\section*{\sffamily\bfseries\color{Deep}The week in one paragraph}
\vspace{-1.5mm}{\color{Teal}\rule{\linewidth}{1.1pt}}\vspace{2.5mm}

\begin{keybox}
{\sffamily\bfseries\small\color{Deep}The biggest regular week this watch has run, and the
measurement side took it --- but the structural story is indoors. Occupational \& indoor
went 7 $\rightarrow$ 19, and what arrived was not exposure epidemiology but
\emph{engineering}: ventilation rates, filter media, purifier siting and smoke evacuation.}\\[1.4mm]
\footnotesize
\textbf{Scale first, because it conditions everything below.} 150 records over seven
issues, the largest week outside the founding backfill, on daily totals of 14, 16, 33, 15,
\textbf{38}, 27 and \textbf{7} --- a 5.4$\times$ range. The 7 on Saturday is not a thin
harvest but a \textbf{true PubMed zero}: no record entering PubMed on 5 September matched
either axis, confirmed independently by the connector against a day that admitted only 544
records in total. Read every week-over-week movement below against that volatility.

\textbf{The indoor turn is the week's real finding.} \textit{Occupational \& indoor} at 19
ties the founding backfill week and is more than double any regular week this watch has
recorded (previous regular maximum 9). The composition of that cluster is what matters:
Jin et al.\ (31 Aug) optimise air-change rate against purifier placement jointly rather
than separately; Wood et al.\ (31 Aug) and Derwein et al.\ (4 Sep) both put ventilation
rate against filtration and mask blocking in the same model; Jamali et al.\ (3 Sep) and
Lv H et al.\ (2 Sep) report new filter media (all-cellulose HEPA, polyimide
microsphere/nanofibrous). Alongside them sit household solid fuel (Gu et al., Chen Y
et al.), surgical smoke (Ma L et al.), welding fume (Yang et al., Wei J et al.), wood dust
(Srinivasan et al.) and chalk dust (Mbazima et al., 5 Sep). \textbf{This is a control- and
mitigation-side literature, not an exposure--response one}, and it is the part of the week
most directly usable by an instrumentation programme.

\textbf{The measurement side is back at its ceiling.} \textit{Exposure assessment \&
modelling} 24 $\rightarrow$ \textbf{47} --- the highest this cluster has ever recorded, and
level with the all-time single-subtopic weekly high set by \textit{Sensing} in the
2--8 August week --- with \textit{Sensing} itself 21 $\rightarrow$ 34. Together 81 of 150,
a measurement-side share of \textbf{54\%}, tying the highest on record.

\textbf{The health side thinned correspondingly, and the effect estimates thinned with
it.} \textit{Respiratory \& allergic} (9 $\rightarrow$ 5) and \textit{Burden, policy \&
mitigation} (13 $\rightarrow$ 7) both hit their lowest weekly counts recorded.
\textbf{34 pooled effect estimates against 150 records} is the weakest estimate density of
any week this watch has run --- one estimate per 4.4 records, against one per 2.7 last
week. \textbf{Operationally the week was clean}: 150 records, \textbf{150 unique DOIs},
so unlike last week there was no double-publication to repair.
\end{keybox}

\clearpage

% -------------------------------------------------------------- trend
\section*{\sffamily\bfseries\color{Deep}Subtopic trend and inflection}
\vspace{-1.5mm}{\color{Teal}\rule{\linewidth}{1.1pt}}\vspace{2.5mm}

\noindent\includegraphics[width=\linewidth]{\FIGDIR/w1_subtopic_trend.png}
\figcap{Per-issue subtopic counts across the week from \texttt{state/metrics.csv}.
\textbf{Read the denominators first}: daily totals were 14, 16, 33, 15, \textbf{38}, 27
and \textbf{7}. The 3 September peak of 38 and the 5 September floor of 7 are both
harvest-side artefacts rather than field signals --- the first a heavy Crossref-by-ISSN
deposit day, the second a weekend on which PubMed admitted no PM record at all. The
Saturday point should not be read as a decline.}

\footnotesize
\textbf{Expanding.} \textit{Occupational \& indoor} \textbf{7 $\rightarrow$ 19} is the
week's largest movement in absolute and relative terms, and its 19 ties the founding
backfill week against a regular-week maximum of 9. \textit{Exposure assessment \&
modelling} \textbf{24 $\rightarrow$ 47} sets a cluster record; \textit{Sensing,
forecasting \& instrumentation} \textbf{21 $\rightarrow$ 34} recovers the ground it lost
in the two preceding weeks. \textit{Reproductive \& developmental}
(\textbf{3 $\rightarrow$ 7}) is at its highest regular-week level, and
\textit{Cardiovascular \& metabolic} (\textbf{5 $\rightarrow$ 7}) continues the slow
recovery the 29 August rollup flagged.

\vspace{1.2mm}
\textbf{Contracting.} \textit{Burden, policy \& mitigation} \textbf{13 $\rightarrow$ 7}
and \textit{Respiratory \& allergic} \textbf{9 $\rightarrow$ 5} are both the lowest weekly
readings this watch has recorded for those clusters. \textit{Neuro / mental health} fell
\textbf{7 $\rightarrow$ 4}; that is low but not a record (the 2--8 August week returned 2).
\textit{Mechanistic toxicology} was effectively flat at \textbf{12 $\rightarrow$ 11},
holding a three-week band of 11--12.

\vspace{1.2mm}
\textbf{The inflection worth naming.} For four weeks the measurement-side share has run
41--54\% with no trend, and the 29 August rollup said explicitly that no claim about the
field should rest on it. That still holds. What is \emph{not} noise is the indoor cluster:
a 7 $\rightarrow$ 19 jump concentrated in ventilation, filtration and source-capture
engineering, arriving in the same week that \textit{Burden, policy \& mitigation} hit its
floor. The mitigation literature did not shrink --- it moved from the policy register into
the engineering register. \textbf{Worth a second week of watching before it is called a
trend}, because a single MDPI or Elsevier volume release can manufacture a cluster of this
size, and the 27 August precedent in the last rollup is exactly that failure.

%% LAYOUT RULE (2026-08-02): architecture donut in the whitespace here, not on its own page.
\vspace{2mm}
\noindent\includegraphics[width=\linewidth]{\FIGDIR/f2a_architecture.png}
\figcap{Study architecture over the pooled week. \textit{Modelling / inventory} (32) leads
\textit{Measurement campaign} (30) --- the first week this watch has recorded modelling
ahead of measurement --- with \textit{Experimental / toxicology} at 18 and
\textit{Metadata only} at 16. Primary observational epidemiology (cohort, cross-sectional,
acute) totals 28, fewer than modelling alone. Only \textbf{one} interventional record
appears in the entire 150.}

\clearpage

% -------------------------------------------------------------- pooled analytics
\section*{\sffamily\bfseries\color{Deep}Pooled analytics}
\vspace{-1.5mm}{\color{Teal}\rule{\linewidth}{1.1pt}}\vspace{3mm}

\noindent\includegraphics[width=\linewidth]{\FIGDIR/f1_subtopics.png}
\figcap{Pooled subtopic distribution, 150 records. All ten clusters were active, for the
sixth consecutive week. The two measurement clusters alone (47 + 34 = 81) outweigh the
eight health and policy clusters combined (69).}

\vspace{2mm}
\noindent\includegraphics[width=\linewidth]{\FIGDIR/f2b_endpoint.png}
\figcap{Health endpoint by organ system. The three largest bars carry \emph{no} human
endpoint --- instrument (16), aerosol chemistry (15) and exposure (14) --- and with
monitoring and modelling added, \textbf{57 of 150 records have no health endpoint at all}.
Respiratory (8), carcinogenic risk (7) and cardiovascular (6) lead the endpoints that
exist. Note that \textit{Cancer} and \textit{Oncologic} remain split in
\texttt{ENDPOINT\_CANON}, so oncological work is understated by this panel; logged
3 September, still unfixed.}

\vspace{2mm}
\noindent\includegraphics[width=\linewidth]{\FIGDIR/f3_metric_geography.png}
\figcap{Left --- particle metric. Unspeciated PM and emission factors (57) lead
PM$_{2.5}$-only (43), with PM$_{2.5}$+PM$_{10}$ jointly at 17. Ultrafine remains
marginal at 4 despite two of the week's strongest records being ultrafine-resolved.
Right --- geographic provenance. \textbf{China supplies 46 of 150 (31\%)}. The panel
understates concentration in the anglophone west because the raw geography field is still
unnormalised --- \textit{North America} (11), \textit{United States} (7) and
\textit{Canada} (1) are distinct labels, as are \textit{Europe} (17),
\textit{United Kingdom} (6), \textit{Germany} (3), \textit{Belgium} (2) and others. That
defect is carried from daily-issue authoring and is logged again below.}

\clearpage

\noindent\includegraphics[width=\linewidth]{\FIGDIR/f5_forest.png}
\vspace{1mm}
\noindent\includegraphics[width=\linewidth]{\FIGDIR/f5_forest_2.png}
\figcap{All 34 pooled effect estimates on a common logarithmic axis.
\textbf{These are not mutually comparable and are not meta-analysed}: contrasts include
per-IQR, per-1\,$\mu$g/m$^3$, per-quartile and AQI-unit changes, on HR (16), RR (9),
OR (7), aOR (1) and IRR (1). \textbf{Heterogeneity is structural, not statistical}, so no
$I^2$ is reported --- pooling these would be a category error. Two features are readable
anyway. First, \textbf{precision tracks sample size to the point of implausibility}: Chen
Y et al.'s congenital heart disease ORs (2.23, 2.19--2.27) and Chen N et al.'s tightest
VTE estimates (1.407, 1.396--1.417) carry intervals so narrow that they imply either
enormous registries or under-propagated uncertainty, and the abstracts do not let a reader
tell which. Second, \textbf{only two of the 34 cross the null} (Batisse et al.'s
paediatric leukaemia and neuroblastoma HRs) --- a distribution that reflects what abstracts
choose to report at least as much as what the studies found.}

\vspace{2mm}
\noindent\includegraphics[width=\linewidth]{\FIGDIR/f6_lifecourse.png}
\figcap{Life-course windows across the pooled week. Working-age adulthood dominates,
consistent with the occupational and indoor surge described above.}

\vspace{2mm}
\noindent\includegraphics[width=\linewidth]{\FIGDIR/f4_heatmap.png}
\figcap{Subtopic $\times$ architecture, pooled. The standing structural finding survives
another week at $n=150$: the \textit{Sensing} and \textit{Exposure} rows carry almost no
health-endpoint linkage, and \textit{Mechanistic toxicology} remains almost entirely
experimental. Instrumentation and epidemiology are still being written by different
people in different journals.}

\clearpage

% -------------------------------------------------------------- venues + trials
\section*{\sffamily\bfseries\color{Deep}Venues, trial watch, and what the pipeline did}
\vspace{-1.5mm}{\color{Teal}\rule{\linewidth}{1.1pt}}\vspace{2.5mm}

\footnotesize
\textbf{Venues.} \textit{Environ Sci Technol} at \textbf{23 of 150} is the most
concentrated single-venue week this watch has recorded, and it is a harvesting artefact as
much as an editorial one --- ES\&T is on the Crossref by-ISSN list and publishes in daily
batches, so it enters this corpus at full rate while journals reached only through PubMed
enter at PubMed's cadence. \textit{Atmos Chem Phys} (8), \textit{J Environ Sci (China)}
(8) and \textit{Atmosphere} (8) follow, then \textit{Environ Pollut} (7).
\textbf{Research Square supplied 5 preprints}, all labelled as such in the register.

\vspace{1.5mm}
\textbf{Recurring authors.} Surname clustering is dominated by high-frequency Chinese
surnames --- Wang (7), Chen (6), Li (4), Liu (4) --- which is a naming artefact, not a
research-group signal, and this watch has no author-disambiguation field to resolve it.
\textbf{No genuine repeat-group cluster is claimed for this week.} Adding a first-author
affiliation field remains the open item that would make this section meaningful; it has
been on the list since 3 September.

\vspace{1.5mm}
\textbf{Trial watch --- one new interventional registration in the period.}
\textbf{NCT07805824} (first posted 4 September) is a cluster-randomised trial of a
workplace heat and air-pollution mitigation package across \textbf{20 restaurants and 200
oven workers} in Najaf, Iraq, with a delayed-intervention comparator over 24 weeks. Its
primary endpoint is an objective gingival melanin index; its \textbf{sole secondary
endpoint is change in personal PM$_{2.5}$ exposure, measured both gravimetrically and in
real time}. \texttt{analyze\_endpoints} was run: 1 primary, 1 secondary, 0 other.
\textbf{Registry-hygiene flags are recorded on the record and should be read before
citing it}: it is labelled PHASE4 with no drug; \texttt{start\_date} equals
\texttt{primary\_completion\_date} equals 2026-01-01 while status is COMPLETED; and the
public \texttt{detailed\_description} still carries the submitter's own pre-submission
checklist. The endpoint design is usable; the timeline fields are not.
Two further records were \emph{updated} rather than registered in the period ---
\textbf{NCT04153539}, the Fudan busy-road versus traffic-free-park crossover in 69 healthy
young adults, and \textbf{NCT05016271} --- and \textbf{NCT07797413} (Savannah
port-adjacent nutrition feasibility trial, $n=20$) was carried in the August monthly.
\texttt{trials.json} now holds \textbf{26 trials across 11 windows}.

\vspace{1.5mm}
\textbf{A weekend zero in the registry is structural.} ClinicalTrials.gov posted
\textbf{no updates registry-wide} on 5 September, nor on 29 or 30 August, against
973--1{,}001 on each intervening weekday. Every Saturday and Sunday daily issue will
therefore report a trial zero for calendar reasons alone. This was verified this week
across six dates and should stop the weekend zeros being mistaken for an absence of PM
trial activity --- a misreading the daily issues were previously set up to invite.

\vspace{1.5mm}
\textbf{What the harvesting legs did.} PubMed carried five of seven days and returned a
verified zero on the sixth-largest\footnote{Ordering by daily count: 38, 33, 27, 16, 15,
14, 7 --- Saturday 5 September is the smallest.} --- on the \emph{smallest} day.
Crossref-by-ISSN was the week's workhorse and the reason ES\&T dominates the venue table.
\textbf{OpenAlex returned HTTP 429 on every day of the week}, unchanged since 28 July, and
is now best described as a dead leg rather than a degraded one. \textbf{arXiv returned
HTTP 200 every day but its newest submission has been frozen at 2026-08-20 all week}, so
it contributed nothing to any issue in this period --- worth a look at the query before
next week. \textbf{Consensus returned a nil for the seventh consecutive run}, but the
4 September check confirmed why: its hits are the right literature with the wrong dates,
and one of them was already in this corpus. Post-filtering Consensus on Crossref
\texttt{created} inside the window remains the recommended fix and is still not
implemented.

\vspace{2mm}
% -------------------------------------------------------------- gaps
\section*{\sffamily\bfseries\color{Deep}Open gaps}
\vspace{-1.5mm}{\color{Teal}\rule{\linewidth}{1.1pt}}\vspace{2.5mm}

\footnotesize
\begin{itemize}
\item \textbf{Ultrafine remains 4 of 150} while two of the week's strongest records ---
Sawan et al.\ on airport-attributable UFP and Shi S et al.\ on ship-plume number modes ---
are ultrafine-resolved. The metric that carries the source signal is still the metric
almost nobody measures.

\item \textbf{One interventional record in 150.} The evidence base moving into the
engineering register this week is almost entirely observational or bench-scale; the
ventilation and filtration papers optimise against models and chambers, not against
health outcomes in occupied buildings.

\item \textbf{No health endpoint in 57 of 150 records}, and the sensing/exposure rows of
the heatmap remain empty of endpoint linkage for the sixth week. This is now the
longest-standing structural finding in the series and it has not moved.

\item \textbf{Estimate density fell to one per 4.4 records.} Whether that is a real
thinning of quantitative epidemiology or an artefact of a measurement-heavy week is not
decidable from one week; the four-week series is 3.3, 2.7, 2.9, 4.4.

\item \textbf{Pipeline defects carried forward}: unnormalised geography labels (which
understate western concentration in f3), the \texttt{Cancer}/\texttt{Oncologic} split in
\texttt{ENDPOINT\_CANON}, no first-author affiliation field, and a 6-element
\texttt{LIFECOURSE} array in the 2026-08-03 corpus where the schema requires 5 --- outside
this window, but it will corrupt any rollup that spans it.
\end{itemize}

\clearpage

% -------------------------------------------------------------- digest
\section*{\sffamily\bfseries\color{Deep}Paper-level digest --- the week by subtopic}
\vspace{-1.5mm}{\color{Teal}\rule{\linewidth}{1.1pt}}\vspace{1mm}

{\fontsize{7.4}{9}\selectfont\color{Slate}Every record that appeared in a daily issue
keeps that issue's critical summary verbatim; nothing here is re-written or re-assessed.
Records surfaced without an abstract are shown as such. Ordered by subtopic, then by the
date of the issue that first carried the record.\par}

\input{digest_body}

% -------------------------------------------------------------- register
\section*{\sffamily\bfseries\color{Deep}Pooled corpus register}
\vspace{-1.5mm}{\color{Teal}\rule{\linewidth}{1.1pt}}\vspace{2.5mm}

{\fontsize{6.7}{8.4}\selectfont
\setlength{\tabcolsep}{2.4pt}
\begin{longtable}{@{}>{\RaggedRight}p{25mm}>{\RaggedRight}p{41mm}>{\RaggedRight}p{27mm}%
>{\RaggedRight}p{23mm}>{\RaggedRight}p{18mm}>{\RaggedRight}p{15mm}@{}}
\toprule
\textsf{\textbf{First author}} & \textsf{\textbf{Focus}} & \textsf{\textbf{Journal}} &
\textsf{\textbf{Design}} & \textsf{\textbf{Metric}} & \textsf{\textbf{Region}}\\
\midrule
\endfirsthead
\toprule
\textsf{\textbf{First author}} & \textsf{\textbf{Focus}} & \textsf{\textbf{Journal}} &
\textsf{\textbf{Design}} & \textsf{\textbf{Metric}} & \textsf{\textbf{Region}}\\
\midrule
\endhead
\multicolumn{6}{@{}l@{}}{\cellcolor{Mist}\textsf{\textbf{\textcolor{Deep}{Neuro / mental health}}}}\\[0.6mm]
\href{https://doi.org/10.1001/jamanetworkopen.2026.25015}{Dales} & Wildfire-sourced PM$_{2.5}$ and ED visits for migraine and other primary headache & \textit{JAMA Netw Open} & Case-crossover & PM$_{2.5}$ (wildfire) & Canada\\
\href{https://doi.org/10.1093/ageing/afag215}{Glover} & Umbrella meta-review of urban environmental exposures and cognitive health & \textit{Age Ageing} & Umbrella review / meta-analysis & PM$_{2.5}$ & Global\\
\href{https://doi.org/10.1016/j.nsa.2026.107019}{Goossens} & In utero air pollution exposure and pre-/postnatal brain development & \textit{Neurosci Appl} & Narrative review & PM$_{2.5}$ / PM$_{10}$ / PM$_{\mathrm{coarse}}$ & Global\\
\href{https://doi.org/10.1016/j.isci.2026.116617}{Ma} & PM$_{2.5}$-induced neurotoxicity and depression: exposure risks and mechanisms & \textit{iScience} & Narrative review & PM$_{2.5}$ + components & Global\\
\href{https://doi.org/10.3760/cma.j.cn112150-20250804-00752}{Pan} & Maternal PM exposure and offspring neurodevelopment: effects and mechanisms & \textit{Zhonghua Yu Fang Yi Xue Za Zhi} & Narrative review & PM (bulk) & China\\
\addlinespace[0.9mm]
\multicolumn{6}{@{}l@{}}{\cellcolor{Mist}\textsf{\textbf{\textcolor{Teal}{Cardiovascular \& metabolic}}}}\\[0.6mm]
\href{https://doi.org/10.3760/cma.j.cn112338-20251120-00843}{Huang} & PM$_{2.5}$ and dyslipidaemia in rural Chinese children; dietary effect modification & \textit{Zhonghua Liu Xing Bing Xue Za Zhi} & Cross-sectional (multistage) & PM$_{2.5}$ & China\\
\href{https://doi.org/10.32394/pe/218966}{Kazmierska} & Air pollution and cardiovascular disease: review of current data & \textit{Przegl Epidemiol} & Narrative review & PM$_{2.5}$ / PM$_{10}$ & Global\\
\href{https://doi.org/10.1007/s00125-026-06814-2}{Xi} & Air pollution, built environment and temperature vs. prediabetes phenotypes & \textit{Diabetologia} & Prospective cohort & PM$_{2.5}$ & Germany\\
\href{https://doi.org/10.1161/JAHA.125.046917}{Yang} & Short-term PM$_{2.5}$ and hospitalisation/ED use in atrial fibrillation & \textit{J Am Heart Assoc} & Case-crossover & PM$_{2.5}$ (wildfire-influenced) & USA\\
\addlinespace[0.9mm]
\multicolumn{6}{@{}l@{}}{\cellcolor{Mist}\textsf{\textbf{\textcolor{Sky}{Reproductive \& developmental}}}}\\[0.6mm]
\href{https://doi.org/10.1016/j.jes.2025.10.043}{Chen} & Prenatal air pollution, fetal growth trajectory phenotypes and child neurodevelopment & \textit{J Environ Sci (China)} & Prospective cohort & PM$_{2.5}$ / PM$_{10}$ & China\\
\href{https://doi.org/10.1016/j.ecoenv.2026.120566}{Huang} & Joint prenatal bisphenol + air-pollutant mixture exposure and preschool IQ & \textit{Ecotoxicol Environ Saf} & Prospective cohort (mixtures) & PM$_{2.5}$ / PM$_{10}$ & China\\
\href{https://doi.org/10.1007/s40572-026-00552-8}{Peterson} & Ambient air pollution, PAHs and pubertal development: critical review & \textit{Curr Environ Health Rep} & Critical review & PM$_{2.5}$ components & Global\\
\href{https://doi.org/10.1016/j.envres.2026.125313}{Xie} & Long-term PM$_{2.5}$ and timing of pubertal development in boys and girls & \textit{Environ Res} & Prospective cohort & PM$_{2.5}$ & China\\
\addlinespace[0.9mm]
\multicolumn{6}{@{}l@{}}{\cellcolor{Mist}\textsf{\textbf{\textcolor{Sage}{Respiratory \& allergic}}}}\\[0.6mm]
\href{https://doi.org/10.1007/s00484-026-03289-z}{Bayar Muluk} & Climate change and otolaryngology: airway disease impacts and emission reduction & \textit{Int J Biometeorol} & Narrative review & PM$_{2.5}$ & Global\\
\href{https://doi.org/10.4103/mgr.MEDGASRES-D-25-00164}{Liao} & Short-term air pollution, respiratory symptoms and lung function in asthma & \textit{Med Gas Res} & Panel study & PM$_{2.5}$ / PM$_{10}$ & China\\
\href{https://doi.org/10.3390/epidemiologia7040097}{Sanchez de Prada} & Ambient PM and asthma-coded admissions in children 0-3 y, Spain & \textit{Epidemiologia} & Ecological panel & PM$_{2.5}$ / PM$_{10}$ & Spain\\
\href{https://doi.org/10.1016/j.jhazmat.2026.143076}{Yang} & Allergenic vs. toxic risk of airborne seed hairs compared with ambient PM$_{2.5}$ & \textit{J Hazard Mater} & Field measurement & PM$_{2.5}$ + bioaerosol & China\\
\addlinespace[0.9mm]
\multicolumn{6}{@{}l@{}}{\cellcolor{Mist}\textsf{\textbf{\textcolor{Violet}{Mechanistic toxicology}}}}\\[0.6mm]
\href{https://doi.org/10.1016/j.jes.2025.12.074}{Ding} & PM$_{2.5}$ promotes vascular calcification via miR-27b-5p/HES1 signalling & \textit{J Environ Sci (China)} & In vitro + organoid + murine & PM$_{2.5}$ & China\\
\href{https://doi.org/10.1007/s00360-026-01698-5}{Maraschi} & Osmoregulatory disruption and metal bioaccumulation from settleable atmospheric PM & \textit{J Comp Physiol B} & Controlled exposure (ecotox) & Settleable APM & Brazil\\
\href{https://doi.org/10.1016/j.jes.2025.11.018}{Wu} & Macrophage NLRP3 inflammasome in real-ambient PM-induced hepatic glucose dysregulation & \textit{J Environ Sci (China)} & Real-ambient murine + co-culture & Real-ambient PM & China\\
\addlinespace[0.9mm]
\multicolumn{6}{@{}l@{}}{\cellcolor{Mist}\textsf{\textbf{\textcolor{Clay}{Occupational \& indoor}}}}\\[0.6mm]
\href{https://doi.org/10.3389/fmedt.2026.1796828}{Chen} & Nested suction instrument to cut intraoperative surgical-smoke PM$_{2.5}$ & \textit{Front Med Technol} & Bench / device experiment & PM$_{2.5}$ (surgical smoke) & China\\
\href{https://doi.org/10.1093/annweh/wxag055}{Cole} & Determinants of respirable crystalline silica exposure in Australian tunnelling & \textit{Ann Work Expo Health} & Observational exposure study & Respirable dust / RCS & Australia\\
\href{https://doi.org/10.1016/j.jhazmat.2026.143081}{Ding} & Stage-resolved hair Zn-isotope dynamics as a firefighting combustion-exposure tracer & \textit{J Hazard Mater} & Repeated-measures biomarker & Combustion-derived PM & China\\
\href{https://doi.org/10.3390/jox16040119}{Ramsden} & Xenobiotic hazards in aircraft cabin air & \textit{J Xenobiot} & Review + risk economics & Ultrafine particles & Global\\
\addlinespace[0.9mm]
\multicolumn{6}{@{}l@{}}{\cellcolor{Mist}\textsf{\textbf{\textcolor{Amber}{Sensing, forecasting \& instrumentation}}}}\\[0.6mm]
\href{https://doi.org/10.1016/j.jenvman.2026.130515}{Shi} & Adaptive spatiotemporal graph transformer for multi-site multi-step PM$_{2.5}$ forecasting & \textit{J Environ Manage} & Deep-learning model & PM$_{2.5}$ & China / India\\
\href{https://doi.org/10.3390/bios16070380}{Tsuchiya} & Seasonal and indoor-outdoor characteristics of fluorescent aerosol particles in offices & \textit{Biosensors} & Real-time sensor campaign & Bioaerosol / coarse & Japan\\
\addlinespace[0.9mm]
\multicolumn{6}{@{}l@{}}{\cellcolor{Mist}\textsf{\textbf{\textcolor{Amber}{Exposure assessment \& modelling}}}}\\[0.6mm]
\href{https://doi.org/10.1029/2025GH001621}{Solorzano-Araque} & Multiscale exposure modelling: LOTOS-EUROS + mobility surveys + particle toxicity & \textit{GeoHealth} & CTM + mobility integration & PM$_{2.5}$ / PM$_{10}$ & Colombia\\
\addlinespace[0.9mm]
\multicolumn{6}{@{}l@{}}{\cellcolor{Mist}\textsf{\textbf{\textcolor{Coral}{Burden, policy \& mitigation}}}}\\[0.6mm]
\href{https://doi.org/10.1002/hsr2.72886}{Karamzad} & Ischaemic heart disease burden in MENA 1990-2021; PM attributable fraction 28.2\% & \textit{Health Sci Rep} & GBD secondary analysis & Ambient PM & MENA\\
\href{https://doi.org/10.1371/journal.pone.0347989}{Rees} & HEPA filtration to reduce respiratory infection in care homes & \textit{PLoS One} & RCT process evaluation & Filtered PM & UK\\
\href{https://doi.org/10.1016/j.jes.2026.01.009}{Shan} & Policy-driven emission reductions and post-pandemic rebound, high-resolution inventories & \textit{J Environ Sci (China)} & Emission inventory & PM emissions & China\\
\addlinespace[0.9mm]
\multicolumn{6}{@{}l@{}}{\cellcolor{Mist}\textsf{\textbf{\textcolor{Slate}{Other clinical endpoints}}}}\\[0.6mm]
\href{https://doi.org/10.3390/dj14070440}{Lupita} & PM$_{2.5}$, preventive-care capacity and caries/edentulism burden across the EU & \textit{Dent J} & Ecological panel & PM$_{2.5}$ & EU-27\\
\href{https://doi.org/10.1016/j.jdin.2026.06.008}{Park} & Long-term PM exposure and risk of androgenetic alopecia & \textit{JAAD Int} & Nationwide cohort & PM$_{2.5}$ / PM$_{10}$ & South Korea\\
\href{https://doi.org/10.3171/2026.3.PEDS2615}{Sun} & Environmental exposure and area deprivation vs. paediatric hydrocephalus outcomes & \textit{J Neurosurg Pediatr} & Registry cohort + geospatial & PM$_{2.5}$ (satellite-derived) & USA\\
\addlinespace[0.9mm]
\bottomrule
\end{longtable}}


\vspace{4mm}
\begin{keybox}
{\sffamily\bfseries\footnotesize\color{Deep}Provenance and method}\\[1.2mm]
{\fontsize{7.2}{8.8}\selectfont
Period \textbf{2026-08-30 to 2026-09-05}, seven daily issues pooled.
\textbf{This rollup queried no API.} It aggregates \texttt{state/corpus/*.json},
\texttt{state/metrics.csv} and \texttt{state/trials.json} over
\texttt{PMRW\_START}--\texttt{PMRW\_END}; every record was screened, summarised and
DOI-gated at daily-issue time and is reproduced verbatim. \textbf{150 records, 150 unique
DOIs} --- no duplicate publication this period, unlike the two repaired last week. Daily
totals: 14, 16, 33, 15, 38, 27, 7. \textbf{Week-over-week comparisons} are against the
23--29 August rollup (109 records), with the 16--22 August (112), 9--15 August (134) and
2--8 August (146) rollups used where a longer run is cited; the 26 July--1 August week
(295) is excluded from "record" claims because it contains the founding retrospective
back-scan. \textbf{Effect estimates} are pooled as recorded, on heterogeneous contrasts,
with no meta-analysis and no heterogeneity statistic. \textbf{Figures} are regenerated
over the pooled range, not composited from the daily images. \textbf{Labels} --- subtopic,
design, particle-metric, geography --- are single-label by dominant contribution, assigned
at daily-issue time, and are a judgement rather than a controlled vocabulary.\par}
\end{keybox}

\end{document}
